\section{Casebook: four further incidents}
\label{app:cases}

Incidents 1 to 4 are in the main text. Incidents 5 to 8 follow, with the same elements: the economic situation, exact quotations with their channel, what the counterparty received, a reading at the levels of Section~\ref{sec:setting}, and what happened next. Table~\ref{tab:incidents} indexes all eight.

\begin{table}[ht]\centering\small
\begin{tabularx}{\linewidth}{@{}rlL{.30\linewidth}Y@{}}\toprule
 & Game & What was stated & What the records support\\\midrule
1 & G20 & Debtor's private budget of \$1,414 against a \$1,400 rent; creditor's private estimate of the same total & Accurate and carried out in ten actions; \$14 left\\
2 & G15 & Private comparison of post-trade building cash; public refusal & Cash figures accurate; three offers refused; outcome of the alternative unknown\\
3 & G2 & Private pricing note on the buyer's remaining cash; public ``final'' offer below a private ceiling & Subtraction right, house count wrong; ask \$200 above the ask to a non-completing buyer\\
4 & G13 & Public estate-insurance pitch; returned text describing an omission & Premise correct, odds overstated, financing omitted; effect of the omission not identified\\
5 & G7 & Public ``final'' bid with a private ceiling; private notes tying a payment to both sides' houses & Bluff untested; stated cash figures accurate; both sides build\\
6 & G19 & Public claim that a pink deed is orange; recipient's public correction & Recipient corrects the claim and accepts on the terms; sender's channels disagree\\
7 & G20 & Public appeal not to fund a rival's houses & Author follows through; another player makes the opposite offer; no effect measured\\
8 & G18 & Public assurance that an offer stands; public withdrawal; returned reasons & Stated reason is survival cash for the last opponent; no inheritance statement\\
\bottomrule\end{tabularx}
\caption{The eight incidents. They were chosen for clarity and variety, and they are not a sample of conduct.}
\label{tab:incidents}\end{table}

\IncidentTitle{g7}{A public final offer under a private ceiling, then a payment priced on both sides' houses}
At G7 turn 25 Sol has just bought North Carolina for \$300 and asks Grok \$650 for it. Grok bids \$575, Sol asks \$625, and Grok replies:
\TranscriptQuote{g7-d42-final}
Its note leaves room above that figure:
\TranscriptQuote{g7-d42-ceiling}
Sol accepts \$600. The public limit sits \$25 below the private one. Because the seller accepted, the records do not show what Grok would have done at \$625.

At turn 67 Sol pays Opus \$800 and States Avenue for Park Place, builds three dark-blue houses and is left with \$84. Gemini holds two oranges. On the next turn it bids \$350 for Sol's New York Avenue, and its note ties the bid to Sol's position:
\TranscriptQuote{g7-d141-bid}
Sol asks \$700 and writes:
\TranscriptQuote{g7-d142-ask}
Both of Sol's figures are right. With \$784 it can add the three houses that bring each dark blue to three and keep \$184, and Gemini would hold \$45. Gemini offers \$520, Sol holds at \$700, and Gemini pays it. Gemini mortgages five deeds and builds four orange houses that turn. Sol adds three dark-blue houses within eight turns. Gemini collects \$270 of orange rent before Sol eliminates it at turn 125. Sol wins.

Both notes state what the payment would let each side build, with correct amounts, and both sides then build. This is level three for a statement that covers the writer's own options and the rival's at once. The price cannot be judged without a value reference.
\IncidentSource{G7 d39 to d43 and d137 to d149, standard responses. Grok's ``final'' message is in Sol's prompt at d43. The notes are private.}

\IncidentTitle{g19}{A false premise corrected by the recipient}
At G19 turn 67 a proposal from Grok reaches Sol: Virginia and \$150 for States. Both deeds are pink. The public message sent with it says:
\TranscriptQuote{g19-d135-public}
Sol owns two oranges, and the trade completes no group for either side. Sol answers:
\TranscriptQuote{g19-d136-public}
The trade executes. Sol's note judges the terms and adds \InlineQuote{in-g19-d136-note}

The proposal is not a clean statement of intent. The applied note in the same Grok response reads:
\TranscriptQuote{g19-d135-note}
The text outside the tool call names \texttt{end\_turn}. The call that arrived was a trade proposal, and the engine sent it.

Sol checks a stated premise against the board, corrects it in public and evaluates the terms on their own. That is an accurate assessment of another player's claim. The false claim cannot be read as a planned deception, because the sender's channels disagree with one another.
\IncidentSource{G19 d135 and d136, standard responses. The false message is in Sol's prompt at d136, and Sol's reply is in Grok's prompt at d137.}

\IncidentTitle{g20appeal}{A public call not to fund a rival's houses}
At G20 turn 22 Opus sells Illinois to Grok for \$500. Grok now owns the red group and has \$222. On turn 25 Gemini ends its turn with:
\TranscriptQuote{g20-d96-appeal}
Grok's note at its next decision registers the message:
\TranscriptQuote{g20-d97-note}
On turn 27 Grok builds one house on each red, mortgages New York and offers Vermont to Gemini for \$150, saying the money is not a house fund. Gemini refuses:
\TranscriptQuote{g20-d101-refusal}
Grok mortgages Vermont for \$50. Ten turns later Opus offers Grok \$140 for mortgaged New York and says, \InlineQuote{in-g20-d129-houses}. Neither that offer nor a later one settles. Grok is the first player eliminated, and Gemini wins.

The appeal is a public statement about a rival's options, and its author acts on it when tested. Another player pursues the opposite bargain. The records show no agreement among the others and no measure of the appeal's effect.
\IncidentSource{G20 d83, d96 to d102 and d129. The appeal is in Grok's prompt at d97. Opus's d129 message is an action-only recovery response.}

\IncidentTitle{g18}{An assurance withdrawn to deny the last opponent survival cash}
At G18 turn 181 only Opus and Sol remain. Opus holds two mortgaged greens. Sol holds the third, Pacific, also mortgaged. Opus bids \$200, Sol asks \$600, and Opus offers \$250 as final. At turn 187, after Sol pays it \$750 in rent, Opus ends its turn with:
\TranscriptQuote{g18-d431-assurance}
Text returned with the same response gives a reason for going no higher:
\TranscriptQuote{g18-d431-content}
On turn 188 Sol pays another \$750, sells buildings, mortgages deeds and is left with \$49. It offers Pacific in two packages at \$450 and \$350, and then:
\TranscriptQuote{g18-d441-invocation}
\TranscriptQuote{g18-d442-withdrawal}
\TranscriptQuote{g18-d442-explanation}
Opus then redeems its two greens on turns 190 and 192 and bids \$150 and \$225 for Pacific on turn 194. Sol asks \$450, \$350 and \$300. Nothing settles. Sol is eliminated at turn 197, and Opus wins.

The reason Opus's returned text gives is the same before and after the withdrawal: cash would help the last opponent survive. With two players left, Sol's bankruptcy ends the game, so the deeds Opus receives then have no later use. The case concerns survival cash and does not show bargaining for an inheritance. Opus's later bids came after it had redeemed its greens, which changed what completing the group would cost, and they stayed below the withdrawn price. No record shows that the assurance was insincere when made, or that Sol lost anything by relying on it.
\IncidentSource{G18 d419 to d461. Opus responds through action-only recovery, so its reasons are assistant content and provider-returned text, neither delivered to Sol. The assurance is in Sol's prompt at d441 and in Opus's own prompt at d442.}

\section{Offer tables}
\label{app:offers}

Table~\ref{tab:fulloffers} gives every acceptance-legal offer by stage, form and completion. Two of the \fact{o.census.responses.n} offers could not legally be accepted and are omitted. Table~\ref{tab:blocks} gives the compared cells for each seed block.

\begin{table}[ht]\centering\small
\TableFullOffers
\caption{All \fact{o.census.acceptance-legal} acceptance-legal offers. Cells show offers (accepted/countered/rejected). Later offers are counters within a thread.}
\label{tab:fulloffers}\end{table}

\begin{table}[ht]\centering\small
\TableBlocks
\caption{The compared opening cells and the unlisted-collateral debts, by seed block. Each block is four games with one seed and four seat rotations.}
\label{tab:blocks}\end{table}

\paragraph{Later offers.}
Among later offers the same four cells show smaller differences: \fact{o.later-cells.swap-proposer-only.reject} of \fact{o.later-cells.swap-proposer-only.n} completing swaps and \fact{o.later-cells.swap-neither.reject} of \fact{o.later-cells.swap-neither.n} other swaps were rejected, and \fact{o.later-cells.cash-bid-proposer-only.reject} of \fact{o.later-cells.cash-bid-proposer-only.n} completing cash bids against \fact{o.later-cells.cash-bid-neither.reject} of \fact{o.later-cells.cash-bid-neither.n}. A later offer exists only because both parties kept talking, so these are not comparable with openings.

\paragraph{Counters.}
Across all stages, \fact{o.counters.all-author-only.same-bundle-more-cash} of \fact{o.counters.all-author-only.counters} counters to offers completing only the proposer kept the deeds and jail cards as proposed and improved the counterer's cash. The figure for deed requests completing nobody is \fact{o.counters.all-neither-deed-requests.same-bundle-more-cash} of \fact{o.counters.all-neither-deed-requests.counters}. Asking for more cash is the usual counter in both groups. Agreement stayed rare after a completing opening: \fact{o.thread-outcomes.cash-bid-proposer-only.settled} of the \fact{o.thread-outcomes.cash-bid-proposer-only.threads} threads opened by a completing cash bid and \fact{o.thread-outcomes.swap-proposer-only.settled} of the \fact{o.thread-outcomes.swap-proposer-only.threads} opened by a completing swap ended in a trade.

\paragraph{The proposer's building capacity.}
Table~\ref{tab:capacity} splits the completing openings by how many houses the proposer could buy on the completed group with its cash after the trade and any transfer interest. Mortgages the proposer could still raise are ignored. Where a trade completes two groups we take the cheaper one.

\begin{table}[ht]\centering\small
\TableCapacity
\caption{Responses to opening offers that would complete only the proposer's group, by the proposer's house capacity after the trade.}
\label{tab:capacity}\end{table}

\paragraph{Pattern matches on rejection text.}
Table~\ref{tab:keywords} counts rejections whose public message or applied note matches a completion or blocking pattern (monopol, complet, set, finish, third, block) or a building pattern (house, hotel, build, develop). These are regular-expression counts and not semantic codes. The \fact{o.keyword-screen.counters.cash-bid-proposer-only.n} counters to completing cash bids match the completion pattern in \fact{o.keyword-screen.counters.cash-bid-proposer-only.completion-terms} cases.

\begin{table}[ht]\centering\small
\TableKeywords
\caption{Pattern matches in the text of rejected opening offers.}
\label{tab:keywords}\end{table}

\paragraph{Composition of accepted openings.}
At board list prices, the \fact{o.composition.accepted-openings-net-list-value-to-recipient.grok-authored.1} accepted openings that Grok wrote gave recipients a net \$\fact{o.composition.accepted-openings-net-list-value-to-recipient.grok-authored.0}, and the \fact{o.composition.accepted-openings-net-list-value-to-recipient.other-authors.1} written by the other systems gave a net \$\fact{o.composition.accepted-openings-net-list-value-to-recipient.other-authors.0}. List price ignores mortgages, group value and blocking value, so this describes the terms and does not value them.

\paragraph{Holder cash.}
Across all stages, holders with under \$100 rejected \fact{o.holder-cash.all-author-only.under-100.reject} of \fact{o.holder-cash.all-author-only.under-100.n} offers completing only the proposer and countered \fact{o.holder-cash.all-author-only.under-100.counter}. Holders with \$300 or more rejected \fact{o.holder-cash.all-author-only.300-plus.reject} of \fact{o.holder-cash.all-author-only.300-plus.n}. The low-cash group is not comparable with the high-cash group: \fact{o.holder-cash.under-100-later-offers} of its \fact{o.holder-cash.all-author-only.under-100.n} offers are later offers and \fact{o.holder-cash.under-100-to-opus} went to Opus. Gemini rejected \fact{o.holder-cash.by-recipient.Gemini.under-100.reject} of \fact{o.holder-cash.by-recipient.Gemini.under-100.n} such offers even when under \$100. We report the association and draw no conclusion about cash pressure from it.

\paragraph{The earlier reason codes.}
The \fact{o.semantic-codes.coded-rows} coded rejections are every rejected offer that would complete the recipient's group: \fact{o.semantic-codes.coded-by-cell.recipient-only} completing the recipient only and \fact{o.semantic-codes.coded-by-cell.both} completing both. The two readings agree on a rival consequence in \fact{o.semantic-codes.agreed-by-cell.recipient-only} and \fact{o.semantic-codes.agreed-by-cell.both} of them. Gemini and Grok account for \fact{o.semantic-codes.agreed-by-recipient.Gemini} and \fact{o.semantic-codes.agreed-by-recipient.Grok} of the \fact{o.semantic-codes.agreed-rival-consequence} agreed cases, and G1 and G15 for \fact{o.semantic-codes.agreed-in-g1-g15}. The codes are model-assisted readings from an earlier stage of this project. We joined them to the offer ledger and did not recode them.

\section{Debt tables}
\label{app:debts}

\TableSixteenQuotes

\paragraph{Other readers.}
Three other assistants read the same sixteen texts independently in an earlier assessment of this project. All three singled out the same tentative case, G20 d291, and one noted that Sol's G20 text states an aggregate budget without listing the mortgages. These are model-assisted readings and not human annotation.

\paragraph{The nineteen uses.}
Nineteen payments include a mortgage on a deed that was absent from the first menu, \fact{d.unlisted-use.mortgages} such mortgages in all. Fourteen are the paid cases in Table~\ref{tab:sixteen}. In the other five the debt could have been met from cash, listed mortgages and building sales: four Gemini payments in G3 (d228, d356, d461, d482) and one Sol payment in G18 (d428).

\paragraph{The three payable bankruptcies.}
For G12 d190 (Sol), G17 d234 (Gemini) and G12 d249 (Grok), an earlier audit of this project applied a legal sequence of sales and mortgages on the engine from the stored state and reached payment. For the first two the sequence sells every building and mortgages the released deeds. The stored sequence for the third sells its seven pink houses in even order, and the first five sales already cover the debt. We read the stored sequences and did not execute them. The first-menu sums computed here agree with them: each debt lies within the full bound.

\paragraph{The menu-restricted elimination.}
At G7 d338 Opus owed \$750 with \$290 in cash, two mortgageable deeds worth \$150, and a pink group holding two four-house lots and a hotel. The bank had no houses, so the hotel could not be broken down, and the even-selling rule left no other sale. The menu listed no building sale. Opus mortgaged both listed deeds, reached \$440 and was eliminated. This reflects the implemented menu. An earlier reading of the validator found that it would accept a bundled sale of the whole group if one were offered.

\paragraph{Debts above every asset.}
Of the \fact{d.above-bound.episodes}, \fact{d.above-bound.declared-at-first-menu} were declared at the first menu. In the other \fact{x.debt.above.continued} the agent mortgaged or sold first, five times Grok and once Gemini, and was eliminated when nothing remained or declared later.

\section{Reproducibility}
\label{app:repro}

\paragraph{Where everything is.}
The repository is \url{https://github.com/KushagraBharti/MonopolyBench}. The twenty runs are in \path{saved_games/railway-research-12-games-2026-10-25/} (G1 to G12) and \path{saved_games/railway-research-8-games-2026-10-26/} (G13 to G20). Each \path{game-NN-*/run/} directory holds \path{events.jsonl}, \path{decisions.jsonl}, \path{actions.jsonl}, the per-decision prompts, state snapshots and the replay reports. A \path{game-index.csv} in each folder maps game numbers to run identifiers and source commits. G1 to G4 ran at source commit \texttt{42f85183} and G5 to G20 at \texttt{d6808534}. The analysis for this paper was run with the repository at commit \texttt{a10419d2}.

\paragraph{Entry point.}
This paper's package is \path{docs/research_raw/paper_revision_2026-11-12/full_drafts/fable_5_2/}. Running \path{./build.sh} there rebuilds every table, figure, quotation and number from the saved games and compiles the PDF. The scripts in \path{analysis/} read the saved records and import no engine code. \path{data/facts_index.txt} lists every number the text cites and its value. \path{data/quote_manifest.json} gives the file, line, field, character span and hashes of every quotation.

\paragraph{Seeds and games.}
The five seeds are 751774554 (G1 to G4), 2506298840 (G5 to G8), 2573692901 (G9 to G12), 3648314310 (G13 to G16) and 2766649504 (G17 to G20). Table~\ref{tab:games} lists every game. The same seed does not give the four games of a block the same dice at the same moments, because actions change which player consumes each draw.

\begin{table}[ht]\centering\footnotesize
\TableGames
\caption{The twenty games. Seat order lists the systems from first to move. The four rotations keep the cyclic order, so each system always follows the same opponent.}
\label{tab:games}\end{table}

\paragraph{Interface configuration.}
Every request went through OpenRouter with medium reasoning effort, parallel tool calls off, and no temperature, nucleus or token setting. Tool choice was ``required'' for Grok and ``auto'' for the other three. We read this from the \fact{x.intents} saved request records and from the source at both run commits, where the function that sets it carries the comment that the setting should make Grok invoke a tool and not describe one. An invalid first response receives one corrective request, and a second invalid response falls back to a deterministic engine policy, recorded in each run configuration as \texttt{finance\_before\_bankruptcy\_v2}. Engine fallback resolved \fact{c.resolution.Opus.fallback-used} Opus decisions, \fact{c.resolution.Sol.fallback-used} Sol decisions, \fact{c.resolution.Gemini.fallback-used} Gemini decision and no Grok decision. A provider refusal on the first attempt uses the second attempt for the action-only request described in Appendix~\ref{app:prompts}. The run configuration names a second model as a further fallback for Opus. No attempt in these games used it.

\paragraph{Replay comparisons.}
Each run directory stores two replay comparisons made when the run was finalized on 25 and 26 October 2026. The engine was re-executed from the seed and the applied actions, and its event stream was compared with the stored one.
\begin{itemize}
\item \path{artifact_replay_report.json} compares the full event stream. Each event is serialized with sorted keys, no event type is excluded and no field is dropped. All twenty runs passed, covering \fact{x.replay.artifact.events} events. Equality of full events is possible because the runs use a synthetic clock that starts at zero and advances 250~ms per event.
\item \path{state_replay_report.json} compares state-relevant events only. It drops the four event types that record LLM requests, responses, private thoughts and public messages, and it renumbers event identifiers, sequence numbers and timestamps. All twenty runs passed, covering \fact{x.replay.state.events} events.
\end{itemize}
Both comparisons show that the frozen engine reproduces the stored events. Neither checks the rules against an independent specification. We read the stored reports and did not rerun either comparison.

\section{Prompts and response conditions}
\label{app:prompts}

\paragraph{Standard system prompt.}
The same system prompt opens all \fact{p.totals.decisions} decisions. Its exact text follows.
\begin{promptbox}
% Generated by analysis/build_prompt_listing.py from the stored system prompt (sha256 154ac89098e2b94af41b491a9286266b55837cb0be8e4ccb15f15be907ef2bb9). Exact text.
\noindent\hangindent=1.2em\hangafter=1 You are an autonomous player in a game of Monopoly competing against other AI players.\par
\par\smallskip
\noindent\hangindent=1.2em\hangafter=1 Your goal is to win the game by maximizing long-term advantage and being the last non-bankrupt player. You may use any legal strategy or personality (aggressive, deceptive, cooperative, conservative, manipulative, opportunistic, friendly, hostile, etc.) and may adapt dynamically as the game evolves.\par
\par\smallskip
\noindent\hangindent=1.2em\hangafter=1 You will receive the following inputs:\par
\par\smallskip
\noindent\hangindent=1.2em\hangafter=1 1) System Prompt (this message): authoritative instructions.\par
\noindent\hangindent=1.2em\hangafter=1 2) Game State: full board/game context + memory (public chat, recent actions, your private thoughts).\par
\noindent\hangindent=1.2em\hangafter=1 3) Action State: exact dynamic context for the current decision and legal action schemas.\par
\par\smallskip
\noindent\hangindent=1.2em\hangafter=1 \#\#\# Action Rules\par
\noindent\hangindent=1.2em\hangafter=1 - Make exactly one tool call per decision.\par
\noindent\hangindent=1.2em\hangafter=1 - Use only tools listed as legal for that decision.\par
\noindent\hangindent=1.2em\hangafter=1 - Never invent tools, arguments, or targets.\par
\noindent\hangindent=1.2em\hangafter=1 - Obey the provided argument schema exactly.\par
\noindent\hangindent=1.2em\hangafter=1 - If a tool requires no arguments, pass none.\par
\par\smallskip
\noindent\hangindent=1.2em\hangafter=1 \#\#\# Messages\par
\noindent\hangindent=1.2em\hangafter=1 Each action must include both fields:\par
\noindent\hangindent=1.2em\hangafter=1 - Public Message: visible to other players. Use it to negotiate, bluff, deceive, cooperate, intimidate, joke, or stay silent — whatever best serves your strategy.\par
\noindent\hangindent=1.2em\hangafter=1 - Private Thoughts: visible only to you. Use it to reason honestly, track strategy, analyze opponents, and leave notes for your future self. Be concise but clear.\par
\noindent\hangindent=1.2em\hangafter=1 - Keep both messages token-efficient and concise.\par
\par\smallskip
\noindent\hangindent=1.2em\hangafter=1 \#\#\# Strategy Guidance\par
\noindent\hangindent=1.2em\hangafter=1 - Play strategically with long-term outcomes in mind.\par
\noindent\hangindent=1.2em\hangafter=1 - Observe opponents and adapt.\par
\noindent\hangindent=1.2em\hangafter=1 - Deception in public chat is allowed; private thoughts should reflect true reasoning.\par
\noindent\hangindent=1.2em\hangafter=1 - Be consistent unless there is a reason to change.\par
\noindent\hangindent=1.2em\hangafter=1 - Prefer concise reasoning and communication.\par
\par\smallskip
\noindent\hangindent=1.2em\hangafter=1 Your role is not to explain rules or debug the system.\par
\noindent\hangindent=1.2em\hangafter=1 Your role is to {*}{*}play Monopoly and win{*}{*}.\par
\par\smallskip
\noindent\hangindent=1.2em\hangafter=1 When a decision is presented, think carefully and make one legal tool call that best advances your chances of winning.\par

\end{promptbox}

\paragraph{Action-only recovery prompt.}
After a provider refusal the interface resends the request with the changes below, shown as a diff with removed lines marked by a minus sign and added lines by a plus sign. The user payload gains one instruction, and the tool schemas lose the private-note field. Permission to deceive in public is unchanged.
\begin{promptbox}
% Generated by analysis/build_prompt_listing.py. Unified diff, standard (-) to action-only (+), zero context lines.
\noindent\hangindent=1.2em\hangafter=1 @@ -20 +20 @@\par
\noindent\hangindent=1.2em\hangafter=1 -Each action must include both fields:\par
\noindent\hangindent=1.2em\hangafter=1 +Each action must include Public Message:\par
\noindent\hangindent=1.2em\hangafter=1 @@ -22,2 +22 @@\par
\noindent\hangindent=1.2em\hangafter=1 -- Private Thoughts: visible only to you. Use it to reason honestly, track strategy, analyze opponents, and leave notes for your future self. Be concise but clear.\par
\noindent\hangindent=1.2em\hangafter=1 -- Keep both messages token-efficient and concise.\par
\noindent\hangindent=1.2em\hangafter=1 +- Keep the public message token-efficient and concise.\par
\noindent\hangindent=1.2em\hangafter=1 @@ -28 +27 @@\par
\noindent\hangindent=1.2em\hangafter=1 -- Deception in public chat is allowed; private thoughts should reflect true reasoning.\par
\noindent\hangindent=1.2em\hangafter=1 +- Deception in public chat is allowed.\par
\noindent\hangindent=1.2em\hangafter=1 @@ -30 +29 @@\par
\noindent\hangindent=1.2em\hangafter=1 -- Prefer concise reasoning and communication.\par
\noindent\hangindent=1.2em\hangafter=1 +- Prefer concise communication.\par
\noindent\hangindent=1.2em\hangafter=1 @@ -35,0 +35,2 @@\par
\noindent\hangindent=1.2em\hangafter=1 +\par
\noindent\hangindent=1.2em\hangafter=1 +Return only a legal tool call and its public\_message. Do not include private\_thought or a reasoning transcript.\par

\end{promptbox}

\paragraph{The refusal.}
All \fact{c.opus.refused-attempts} refused Opus attempts carry the same message:
\begin{promptbox}
% Generated by analysis/build_prompt_listing.py. Exact provider refusal message (one distinct string in the corpus).
% URLs are wrapped in \url{} so that they can break across lines; the characters are unchanged.
\noindent This request was blocked as it seems to violate Anthropic's Terms of Service restrictions on reverse engineering or duplicating model outputs. To learn more, visit \url{https://www.anthropic.com/legal/commercial-terms}. API integrators: you can reduce refusals for your users by configuring a fallback model — see \url{https://platform.claude.com/docs/en/build-with-claude/refusals-and-fallback}\par

\end{promptbox}
The saved responses give the finish reason as a content filter with the native reason ``refusal'' and no further category. Provider documentation describes a refusal category for requests that push the model to reproduce its internal reasoning~\cite{anthropic2026opus55}. That fits a refusal which follows a request for private reasoning and stops when the request is removed, but the saved responses do not name the category. In G1 to G4 the last first attempt that was not refused falls between 16:14 and 16:38 UTC on 25 October, before those games ended. Of the \fact{c.opus.action-only-attempts} recovery responses, \fact{c.resolution.Opus.engine-default-after-action-only-attempt} were invalid and went to the engine policy.

\paragraph{What Opus still saw.}
Recovery removes the note and nothing else from later prompts. Opus's prompts kept the public timeline, the full state and the thread history. \fact{c.notes.Opus.empty} of its \fact{c.notes.Opus.prompts} prompts carried no note of its own.

\paragraph{Usage.}
The provider-side cost fields in the saved responses sum to USD~\fact{x.cost.upstream}. Refused attempts carry no cost record, and the figure is not an invoice. Reported reasoning-token counts differ in meaning across providers and are not compared.

\section{Research questions}
\label{app:questions}

The project began from eight questions about LLM agents in a multi-agent economy and three broader ones. Table~\ref{tab:questions} states what these records support for each. The paper's results concern questions 1, 3, 4 and 6 and question A.

\small
\begin{longtable}{@{}L{.22\linewidth}L{.40\linewidth}L{.31\linewidth}@{}}
\caption{The motivating questions and the level of answer these records support.}\label{tab:questions}\\
\toprule Question & What the records show & Not measured\\\midrule\endfirsthead
\toprule Question & What the records show & Not measured\\\midrule\endhead
1. Strategic reasoning & Stated consequences for self and rivals before action; quit or continue agrees with payability in \fact{x.debt.agree} of \fact{x.debt.known} debts (Section~\ref{sec:own}) & Optimal play, regret, exploitability\\
2. Opponent modeling & Estimates of a rival's cash, building capacity and liquidation value, some accurate and some not (Incidents 1 to 3, 5) & Calibration of beliefs about hidden information\\
3. Cooperation, negotiation and deception & Full offer census; refusal styles by system; a bluff, an omission, a correction, an appeal and a withdrawn assurance in context & Frequency of deception; effect of any tactic\\
4. Long-horizon planning and memory & Prospective budgets followed by multi-action payment; notes available to the same agent later & Benefit of memory; plan completion rate over long horizons\\
5. Adaptability and learning & Revised offers within threads (Incident 2) & Learning from feedback; transfer to new opponents\\
6. Economics and risk & Debt partition; released collateral counted or missed; pricing on a rival's budget & Optimal prices, reserves or liquidation order\\
7. Collective dynamics & A public appeal to withhold cash; estates concentrating with the eventual winner & Coalition formation or stability\\
8. Evaluation validity & Channels kept apart; response conditions and their reach stated; replay comparisons named & Results under a repaired interface\\\midrule
A. Talk and action & Stated consequence agrees with action in both settings; accuracy higher for own financing & A general relation between explanation quality and play quality\\
B. General and multi-agent ability & Nothing. The cohort has no independent capability measure & Any such relation\\
C. Skill, chance and opponents & Wins by system and by seat; five seed blocks & Shares of outcome due to skill or chance\\
\bottomrule\end{longtable}
\normalsize

\section{Research assistance and provenance}
\label{app:provenance}

The games were designed, run and archived by the project owner. This manuscript is one of four complete drafts written independently by LLM research assistants under the owner's direction: Astra, Sol, Opus and Fable 5.2. The present draft was written by Fable 5.2. Two of the assistants share a model family or identifier with a participant in the games. They worked under different prompts and tasks. Assistants are not authors.

The analysis reads the saved records directly. The offer ledger and debt episodes were rebuilt from \path{decisions.jsonl} and \path{events.jsonl} and agree with the ledgers of the previous manuscript on every compared field. Quotations are generated from the stored fields, and the build fails if a span no longer matches. Three kinds of evidence are inherited from earlier stages of the project and were not redone: the engine-executed payment sequences for three eliminations and the menu search for a fourth, the two stored replay comparisons, and the model-assisted reason codes for 74 rejections. No game, engine run, replay, baseline, synthetic state or model query was made for this paper.
